
\section{Funkcja straty}
\label{sec:strata}

Podstawowym celem uczenia jest poprawne odtworzenie tokenów usuniętych podczas maskowania. Model przewiduje rozkład prawdopodobieństwa dla każdej zamaskowanej pozycji, a następnie porównuje go z prawidłowym tokenem pochodzącym z oryginalnego programu. Podstawową funkcją straty jest entropia krzyżowa, obliczana wyłącznie dla pozycji, które zostały ukryte. Tokeny pozostawione bez zmian nie uczestniczą bezpośrednio w obliczaniu tej części straty.

Każdy błędnie przewidziany token wpływa na wynik zgodnie z prawdopodobieństwem przypisanym poprawnej odpowiedzi. Projekt umożliwia również uczenie z wieloma poprawnymi rozwiązaniami tego samego zadania. Jeżeli dla jednej instrukcji dostępnych jest kilka równoważnych programów, przewidywanie modelu może być porównywane z referencją najlepiej zgodną z widoczną częścią kodu. Dzięki temu model nie jest niepotrzebnie karany za wygenerowanie poprawnego rozwiązania, które różni się od jednego arbitralnie wybranego programu referencyjnego.

Strata strukturalna wykorzystuje abstrakcyjne drzewo składniowe programu. Reprezentacja generowanego kodu jest wówczas porównywana z reprezentacją sekwencji elementów drzewa składniowego. Do ich dopasowania wykorzystywany jest algorytm Soft-DTW, który pozwala porównywać sekwencje o różnych długościach i uwzględnia możliwe przesunięcia pomiędzy elementami.

Porównywana jest też strata dla właściwego opisu zadania oraz dla opisu pochodzącego z innego przykładu. Model jest karany, jeżeli nie wykorzystuje poprawnej instrukcji wyraźnie lepiej niż instrukcji niedopasowanej. Ma to ograniczyć generowanie kodu opartego wyłącznie na częstości wzorców występujących w zbiorze danych.
